% -----------------------------------------------------------------------------
% PART A
% -----------------------------------------------------------------------------
\subsection*{Part A: The Idea}
\textit{These questions test the mathematical definitions and calculations devoid of physical context.}

\begin{enumerate}[label=\textbf{Q\arabic*.}, leftmargin=*]
    
    \item \textbf{Pearson Correlation Coefficient ($r$)} \\
    Consider the following abstract dataset of $n=5$ coordinate pairs: 
    \[ X = \{ 1, 2, 3, 4, 5 \} \]
    \[ Y = \{ 2, 3, 5, 4, 6 \} \]
    \begin{enumerate}[label=\arabic*.]
        \item Calculate the sample means, $\bar{x}$ and $\bar{y}$.
        \item Calculate the Sum of Squares for the spread of $X$ ($S_{xx}$), the spread of $Y$ ($S_{yy}$), and the covariance ($S_{xy}$).
        \item Using your answers from above, calculate the Pearson correlation coefficient ($r$). Is this a strong or weak correlation?
    \end{enumerate}
    \vspace{0.3cm}

    \item \textbf{Simple Linear Regression} \\
    Using the exact same dataset and the Sum of Squares values calculated in Q1:
    \begin{enumerate}[label=\arabic*.]
        \item Calculate the gradient ($m$) of the least-squares line of best fit.
        \item Calculate the y-intercept ($c$). 
        \item Write down the final equation of the line in the form $\hat{y} = mx + c$.
        \item Calculate the Coefficient of Determination ($R^2$). 
    \end{enumerate}
    \vspace{0.3cm}

    \item \textbf{Linearisation (Mapping)} \\
    A mathematical model follows an exponential curve given by the equation:
    \[ y = a e^{bx} \]
    \begin{enumerate}[label=\arabic*.]
        \item By applying the natural logarithm ($\ln$) to both sides, map this equation into a linear format ($\hat{Y} = mX + C$).
        \item In this new linear space, what does the gradient ($m$) represent in terms of the original constants?
        \item What does the y-intercept ($C$) represent in terms of the original constants?
    \end{enumerate}

\end{enumerate}

\vspace{0.8cm}

% -----------------------------------------------------------------------------
% PART B
% -----------------------------------------------------------------------------
\subsection*{Part B: Exam style questions: Engineering Application}
\textit{These questions test the exact same mathematical concepts as Part A, but applied to the Electrical Engineering scenarios discussed in the lecture.}

\begin{enumerate}[label=\textbf{Q\arabic*.}, leftmargin=*, start=4]

    \item \textbf{Sensor Calibration \& Goodness of Fit} \\
    An engineer is calibrating a pressure sensor. They apply known forces ($X$, in Newtons) and measure the output voltage ($Y$, in millivolts). Rather than providing the raw data, the calibration software outputs the following summary statistics for $n=10$ tests:
    \[ \bar{x} = 20\text{ N}, \quad \bar{y} = 45\text{ mV} \]
    \[ S_{xx} = 400, \quad S_{yy} = 100, \quad S_{xy} = 190 \]
    \begin{enumerate}[label=\arabic*.]
        \item Calculate the Pearson correlation coefficient ($r$) for this sensor.
        \item Calculate the Coefficient of Determination ($R^2$).
        \item \textit{Engineering Interpretation:} Explain exactly what this $R^2$ value tells the engineer about the reliability of the pressure sensor's voltage output. 
    \end{enumerate}
    \vspace{0.3cm}

    \item \textbf{Predictive Modelling \& Tolerances} \\
    Using the summary statistics from Q4, the engineer wants to build the linear regression model $\hat{y} = mx + c$ so the microcontroller can predict the voltage automatically.
    \begin{enumerate}[label=\arabic*.]
        \item Calculate the slope ($m$) and intercept ($c$) of the regression line.
        \item \textit{Terminology:} What is the physical meaning of the gradient ($m$) in this context? (Hint: Use the term \textit{Sensitivity}).
        \item \textit{Terminology:} What is the physical meaning of the intercept ($c$) in this context? (Hint: Think about what happens when $0\text{ N}$ of force is applied).
        \item If a force of $25\text{ N}$ is applied to the sensor, what voltage drop should the system expect?
    \end{enumerate}
    \vspace{0.3cm}

    \item \textbf{Non-Linearity in Circuits (Capacitor Discharge)} \\
    You are measuring the voltage ($V$) across a discharging capacitor over time ($t$, in seconds). The physical law governing this is intrinsically non-linear:
    \[ V = V_0 e^{-\frac{t}{\tau}} \]
    Where $V_0$ is the initial voltage and $\tau$ is the RC time constant of the circuit. 
    
    You have a raw dataset of times and voltages. You cannot fit a straight line to this raw data.
    \begin{enumerate}[label=\arabic*.]
        \item How must you transform your raw data table columns (Time and Voltage) before running a linear regression analysis in Python (\texttt{scipy.stats})?
        \item After transforming the data and calculating the line of best fit, Python tells you the gradient (slope) of the line is $-0.2$. What is the RC time constant ($\tau$) of this circuit?
    \end{enumerate}
    \vspace{0.3cm}

    \item \textbf{Correlation vs. Causation} \\
    A data centre manager notices a strong positive correlation ($r = 0.85$) between the speed of the server cooling fans (RPM) and the number of network packet routing errors occurring per minute. They conclude that the vibration from the high-speed fans is causing the network cables to drop packets, and propose limiting the maximum fan speed to fix the network.
    \begin{enumerate}[label=\arabic*.]
        \item Identify the hidden \lq lurking variable\rq\ that is actually causing both the fans to speed up and the network errors to increase simultaneously. 
        \item What catastrophic engineering failure will likely occur if the manager implements their proposed solution?
    \end{enumerate}

\end{enumerate}
